\documentclass[11pt]{article}
\usepackage[a4paper,margin=1in]{geometry}
\usepackage{amsmath,amssymb}
\usepackage{enumitem}
\usepackage{graphicx}
\usepackage[colorlinks=true,linkcolor=blue,urlcolor=blue,citecolor=blue]{hyperref}
\usepackage{titlesec}

% Footnote symbols (*, dagger, ...) instead of numbers, standard for author affiliations
\renewcommand{\thefootnote}{\fnsymbol{footnote}}

% Simple boxed placeholder for figures/diagrams that are not ready yet.
% Usage: \FigurePlaceholder{Caption text describing what will go here}{label}


\begin{document}

\begin{center}



{\LARGE\bfseries Contract Playbook Comparison System}\\[6pt]
{\large A Retrieval-Augmented System for Automated Contract Clause Review}\\[12pt]

{\normalsize Project Report: Prototype Stage, UCS503}\\[4pt]
{\normalsize Submitted to: Dr.\ Jeelani Asif}\\[4pt]
{\normalsize Thapar Institute of Engineering and Technology}\\[18pt]

\begin{figure}[h]
    \centering
    \includegraphics[width=0.4\textwidth]{image (1).png}
    
\end{figure}
\begin{tabular}{ccc}
\textbf{Mannat, AIML} & \textbf{Upasna Wadhwa, AIML} & \textbf{Yuvicka, AIML} \\[4pt]
\href{mailto:mmannat_be24@thapar.edu}{mmannat\_be24@thapar.edu} &
\href{mailto:uwadhwa_be24@thapar.edu}{uwadhwa\_be24@thapar.edu} &
\href{mailto:yyuvicka_be24@thapar.edu}{yyuvicka\_be24@thapar.edu} \\[4pt]
Roll No: 1024240023 & Roll No: 1024240022 & Roll No: 1024240016
\end{tabular}

\\[18pt]

{\normalsize September 14, 2026}

\end{center}

\vspace{10pt}

\tableofcontents
\vspace{10pt}

\section{Introduction}

\subsection{Background and Context}

Legal and procurement teams at most companies work against an internal
\emph{playbook} or we can say a set of pre-approved clause positions that represent
acceptable risk, such as minimum liability caps, maximum termination
notice periods, or compulsory auto-renewal opt-outs. Every time a new
vendor, SaaS, or service contract comes in, someone on the team has to
read through it (often 20 to 100+ pages), find each clause that matters,
and check it against the playbook by hand. In practice this is slow, it
does not scale well as the number of contracts grows, and different
reviewers tend to catch different things, so review quality ends up
depending on who did it. On top of that, verdicts are rarely tied back
to the exact piece of text that produced them, which makes them hard to
defend or dispute later if a disagreement comes up.

It is tempting to just hand the whole contract to a general-purpose LLM
and ask ``does this match our policy?'', but this is exactly the kind of
use case where that approach breaks down. An LLM asked to reason over an
entire document without being anchored to specific, verifiable text can
produce answers that sound confident but are not actually grounded in
anything and when a missed liability clause has real financial
consequences, an unverifiable answer is not good enough.

\subsection{Problem Statement}

Given a new contract and a defined playbook of clause positions, we want
a system that can locate the relevant section of the contract for each
playbook clause, compare it against the expected position, and produce a
verdict (Match, Deviation, or Missing) that is backed by an exact
citation into the source contract rather than an LLM's unverified
opinion of the whole document. Contract Lifecycle Management is already
an established, funded software category (Ironclad, LinkSquares,
Lawgeex, Evisort, among others), which tells us this is a real gap in
how contracts get reviewed today, not a problem we invented for the
sake of the project.

\subsection{Project Objectives}

\begin{enumerate}
  \item \textbf{Objective 1:} Build a retrieval-augmented pipeline that,
    for a given playbook clause type, retrieves the specific contract
    chunk(s) relevant to it using hybrid semantic and keyword search,
    rather than relying on an LLM to search the whole document itself.
  \item \textbf{Objective 2:} Use an LLM strictly as a comparison step such that
    it only ever sees the playbook position and the retrieved chunk(s),
    never the full contract to classify each clause as Match,
    Deviation, or Missing with a short, text-grounded explanation.
  \item \textbf{Objective 3:} Produce a structured, citation-backed
    report (one row per playbook clause, with status, explanation, and
    an exact reference into the source contract) through a simple demo
    interface, so the output is something a reviewer could actually use
    and trust.
\end{enumerate}

These objectives are specific and measurable because each stage
(retrieval, comparison, reporting) has its own evaluation criterion,
described in Section~\ref{sec:results}. They are achievable within a
single semester because we deliberately scope down to 5-8 clause
categories and a well-established public dataset (CUAD) instead of
building a production-grade system. They are relevant to a genuine,
recognised industry problem, and they are time-bound to the prototype
and final submission milestones of this course.

\section{Methodology and Design}

\subsection{System Architecture}

The system is organised as a staged retrieval-augmented pipeline, split
into three parts that map onto how our team divided the implementation
work: ingestion and retrieval, comparison and evaluation, and reporting.

\subsubsection{High-Level Design}

At a high level, an incoming contract is parsed and broken into
indexable chunks. For every playbook clause type, the retrieval stage
returns the top-$k$ candidate chunks from the contract using a hybrid of
semantic (embedding-based) and keyword (BM25) search. If nothing clears
a minimum similarity threshold, the clause is marked \emph{Missing}
immediately, the LLM is never asked to compare against nothing, which
is one of our main defences against hallucinated verdicts. Otherwise,
the retrieved text and the playbook's documented position are passed to
an LLM, which classifies the clause as \emph{Match} or \emph{Deviation}
and returns a one-sentence explanation. All of this is finally
assembled into a single structured report: one row per playbook clause,
its status, the explanation, and a citation back into the exact part of
the contract that produced the verdict.


\begin{figure}[h]
    \centering
    \includegraphics[width=0.8\textwidth]{pipeline.png}
    \caption{End-to-end pipeline architecture diagram: contract
  ingestion, embedding/indexing, hybrid clause retrieval, LLM-based
  comparison, structured reporting, and evaluation.}
\end{figure}


\begin{figure}[h]
    \centering
    \includegraphics[width=0.4\textwidth]{flow.png}
    \caption{Core workflow diagram showing the decision logic:
  contract in $\rightarrow$ chunking $\rightarrow$ retrieval
  $\rightarrow$ threshold check $\rightarrow$ Match / Deviation /
  Missing $\rightarrow$ cited report out.}
\end{figure}


\subsubsection{Technical Stack}

\begin{itemize}
  \item \textbf{Language:} Python.
  \item \textbf{Embeddings:} \texttt{sentence-transformers}
    (\texttt{all-MiniLM-L6-v2} in the current prototype).
  \item \textbf{Vector store / keyword index:} FAISS for semantic
    retrieval, paired with a BM25 keyword index (\texttt{rank\_bm25}),
    since exact legal terminology such as ``indemnify'' or ``force
    majeure'' often carries more signal than pure semantic similarity.
  \item \textbf{Comparison / LLM:} an accessible LLM API, used only for
    the strictly-scoped comparison step described above, with a
    rule-based fallback for select clause types.
  \item \textbf{Reporting:} a Streamlit application that renders the
    structured report and supports CSV and PDF export.
  \item \textbf{Dataset:} CUAD (Contract Understanding Atticus
    Dataset), supplemented with a small set of synthetic sample
    contracts for local development and demos.
  \item \textbf{Testing / CI:} \texttt{pytest} for unit tests, run on
    every change.
\end{itemize}

\subsection{Implementation Details}

The implementation is split across three modules, each owned by one
team member and developed on its own feature branch before being
merged into \texttt{master}.

\subsubsection{Module 1: Ingestion and Retrieval}

This module parses contract text, chunks it into segments of roughly
300-500 tokens with about 15\% overlap (tuned so a single legal clause
does not get split awkwardly across chunk boundaries), and builds the
hybrid semantic and keyword index used by retrieval. Given a playbook
clause type, it returns the top-$k$ candidate chunks along with a
similarity score, which the pipeline later uses to decide whether to
attempt a comparison at all or fall back directly to \emph{Missing}.

\subsubsection{Module 2: Comparison and Evaluation}

This module takes the chunks returned by retrieval and the documented
playbook position for that clause type, and prompts the LLM to return a
strict Match / Deviation classification with a short explanation, with
a rule-based fallback path for clause types where a deterministic check
is reliable enough (for example, governing-law and termination
clauses). The prompt is deliberately constrained to only use the
provided text, and to say it cannot determine an answer if the
retrieved text is ambiguous, rather than guessing. This module also
owns the evaluation scripts that will check retrieval and
classification quality against CUAD's ground-truth spans and our own
hand-labelled test set.

\subsubsection{Module 3: Reporting (Streamlit App)}

This module takes the structured output of the comparison stage
as one entry per playbook clause, with a status, an explanation, and a
reference back into the contract and turns it into something a
reviewer can actually read and use. Concretely, it currently provides:

\begin{itemize}
  \item A Streamlit interface (\texttt{app.py}) that lets a user select
    a contract and view its clause-by-clause report as a table, with
    each row's status, explanation, and source citation.
  \item A tabular report builder (\texttt{report\_builder.py}) that
    assembles the raw comparison output into a consistent, presentable
    structure before it is rendered or exported.
  \item CSV export (\texttt{csv\_report.py}) for reviewers who want to
    work with the results in a spreadsheet.
  \item PDF export (\texttt{pdf\_report.py}) for a shareable, printable
    version of the same report.
\end{itemize}


\begin{figure}[h]
    \centering
    \includegraphics[width=1.0\textwidth]{Screenshot.png}
    \caption{Screenshot of the Streamlit reporting app, showing a
  generated clause-by-clause report table for a sample contract.}
\end{figure}

\section{Results and Evaluation}
\label{sec:results}

\subsection{Testing Methodology}

\begin{itemize}
  \item \textbf{Unit Testing:} Each module has its own \texttt{pytest}
    test suite. For example, \texttt{test\_ingestion.py} checks
    chunking behaviour and retrieval hit rate, \texttt{test\_comparison.py}
    checks the rule-based fallback logic for specific clause types, and
    \texttt{test\_reporting.py} checks that a tabular report is built
    correctly from comparison output.
  \item \textbf{Integration Testing:} We run the pipeline end to end on
    the synthetic sample contracts included in the repository to
    confirm ingestion, retrieval, comparison, and reporting connect
    correctly as a single flow, rather than only working in isolation.
  \item \textbf{Performance and User Testing:} Not yet carried out at
    the prototype stage. We plan to measure retrieval and end-to-end
    pipeline timing, and to have a teammate outside the reporting
    module walk through the Streamlit app as an informal usability
    check, before final submission.
\end{itemize}

At the time of writing, our full unit test suite passes:

\begin{table}[h]
\centering
\begin{tabular}{|l|l|c|}
  \hline
  \textbf{Test File} & \textbf{Test} & \textbf{Status} \\
  \hline
  test\_ingestion.py & test\_chunking & Passed \\
  \hline
  test\_ingestion.py & test\_retrieval\_hit\_rate & Passed \\
  \hline
  test\_comparison.py & test\_rule\_fallback\_governing\_law & Passed \\
  \hline
  test\_comparison.py & test\_rule\_fallback\_termination & Passed \\
  \hline
  test\_reporting.py & test\_build\_tabular\_report & Passed \\
  \hline
\end{tabular}
\caption{Current unit test suite status (5/5 passing) as of the
  prototype submission.}
\label{tab:unit-tests}
\end{table}

\subsection{Analysis}

At the prototype stage, our main goal was to prove that each stage of
the pipeline works on its own and connects cleanly to the next one,
rather than to hit final accuracy numbers. On that front, we are on
track: all three modules run independently, their unit tests pass
(Table~\ref{tab:unit-tests}), and we can take a sample contract all the
way from raw text to a cited report through the Streamlit app.

We have deliberately not reported retrieval or classification accuracy
numbers (Recall@$k$, MRR, Match/Deviation/Missing agreement) in this
report, since the evaluation scripts for those metrics have not yet
been run against the held-out CUAD split hence reporting numbers before
that would risk giving a misleading picture of where the system
actually stands. That evaluation is the main piece of work planned
before final submission.

The biggest challenge so far has been environment and dependency
management across three people working on separate branches and
getting everyone onto compatible Python and package versions, and
keeping shared module boundaries (in particular, a shared
\texttt{common} configuration module used by all three parts)
consistent so a change in one branch does not silently break another.
We addressed this by testing each feature branch independently before
opening a pull request, and by agreeing on a fixed
\texttt{requirements.txt} early on so dependency versions do not drift
between team members.



\section{Conclusion}

\subsection{Summary of Work}

At this stage, we have a working prototype of the full pipeline:
contract ingestion and chunking, hybrid retrieval against a defined
playbook, LLM-based clause comparison constrained to retrieved text
only (with a rule-based fallback where appropriate), and a
Streamlit-based reporting layer with CSV and PDF export. All three
modules pass their unit tests individually, and we have verified the
pipeline end to end on sample contracts. Of our three objectives, the
first two (retrieval and constrained LLM comparison) are functionally
in place, and the third (a usable, citation-backed report) is
implemented and demoable, though we still need to run it against a
larger, held-out slice of CUAD to get real accuracy numbers rather than
just confirming it works.

\subsection{Key Findings}

\begin{itemize}
  \item \textbf{Finding 1:} Splitting the pipeline into independently
    testable modules (ingestion/retrieval, comparison/evaluation,
    reporting) made it realistic for three people to work in parallel
    on separate branches without constantly blocking each other.
  \item \textbf{Finding 2:} Forcing the LLM comparison step to only see
    retrieved text and to fall back to \emph{Missing} automatically
    when retrieval confidence is low meaningfully reduces the
    opportunity for hallucinated verdicts, compared to asking an LLM to
    reason over a whole contract.
  \item \textbf{Finding 3}:{Hybrid BM25 + semantic retrieval is deferred to Stage 2. Among our four clause categories, Governing Law and Auto-Renewal were hardest to retrieve at rank-1 (43\% and 56\% respectively) versus 90\%+ for the other two, due to query-phrasing and chunk-size mismatches rather than a fundamental retrieval failure; both were resolved to 100\% Recall@3 after targeted fixes.} 
    
\end{itemize}

\subsection{Future Work and Recommendations}

\begin{enumerate}
  \item Run full retrieval and classification evaluation against the
    held-out CUAD split and report Recall@$k$, MRR, and classification
    accuracy per clause category, rather than only confirming the
    pipeline runs end to end.
  \item Complete the hybrid (semantic + BM25) versus semantic-only
    retrieval ablation to justify the hybrid design choice with numbers
    rather than intuition.
  \item Add a faithfulness check (manual, and optionally an NLI-based
    check similar to RAGAS) to verify that LLM explanations are
    actually grounded in the retrieved text, not just
    plausible-sounding.
  \item Expand the Streamlit reporting app based on informal user
    testing feedback .For example, letting a reviewer flag a verdict
    they disagree with, if time allows.
  \item Longer term, consider whether the playbook clause set could be
    extended beyond the initial 5-8 categories, and what would be
    needed to move from a local demo deployment toward something
    closer to production use.
\end{enumerate}

\subsection{Final Remarks}

Building this prototype made the gap between ``an LLM that sounds
right'' and ``an LLM whose answer you can actually check'' much more
concrete than it was when we first scoped the project. Constraining
the comparison step to only retrieved, cited text rather than
letting the model read the whole contract and trust its judgement, forced us to think carefully about failure modes (like what happens
when retrieval finds nothing) rather than assuming the LLM would
handle ambiguity gracefully. We are reasonably happy with where the
prototype stands, and the main work left before final submission is
running the pipeline at scale against CUAD and turning those results
into the quantitative analysis this report currently leaves as a
placeholder.

\end{document}